%% kapa: exam=1 round=34 year=2023 period=1 form=A subject=civil_law subject_ko=민법 seq=24 number=24
%% source: Q-Net 감정평가사 기출문제 (cst00309) artlSeq=5211789 file=raw/1008/5211789/2023년도 제34회 감정평가사 1차시험 1교시.pdf
%% conversion: kordoc 4.24.51 -> tools/gen1.py (auto); review=proofread-sonnet; flags=none
\begin{question}{24}
등기를 마치지 않더라도 물권변동의 효력이 발생하는 경우는? (다툼이 있으면 판례에 따름)
\begin{choices}
  \item 지상권설정계약에 따른 지상권의 취득
  \item 피담보채권의 시효소멸에 따른 저당권의 소멸
  \item 공익사업에 필요한 토지에 관하여 토지소유자와 관계인 사이의 협의에 의한 토지소유권의 취득
  \item 공유토지의 현물분할에 관한 조정조서의 작성에 따른 공유관계의 소멸
  \item 당사자 사이의 법률행위를 원인으로 한 소유권이전등기절차 이행의 소에서의 승소판결에 따른 소유권의 취득
\end{choices}
\end{question}
\answer{2}
